\documentclass[10pt,a4paper]{article}
\usepackage{blindtext}
\usepackage{subcaption}
\usepackage{graphicx}
\usepackage{tikz}
\usepackage{amssymb}
\usepackage{caption}
\usepackage{amsmath}
\usepackage{circuitikz}
\usepackage{hyperref}
\usepackage{amssymb}
\usepackage{amsmath}
\usepackage{listings}

\lstset{
    inputencoding=utf8,
    extendedchars=true,
    literate={á}{{\'a}}1 {é}{{\'e}}1 {í}{{\'i}}1 {ó}{{\'o}}1 {ú}{{\'u}}1 {ñ}{{\~n}}1 {Á}{{\'A}}1 {É}{{\'E}}1 {Í}{{\'I}}1 {Ó}{{\'O}}1 {Ú}{{\'U}}1 {Ñ}{{\~N}}1
}
\usepackage[spanish,activeacute,es-tabla]{babel}
\usepackage[utf8]{inputenc}
\usepackage{ifthen}
\usepackage{listings}
\usepackage{dsfont}
\usepackage{subcaption}
\usepackage{amsmath}
\usepackage[strict]{changepage}
\usepackage[top=1cm,bottom=2cm,left=1cm,right=1cm]{geometry}%
\usepackage{color}%
\newcommand{\tocarEspacios}{%
	\addtolength{\leftskip}{3em}%
	\setlength{\parindent}{0em}%
}

% Especificacion de procs

\newcommand{\In}{\textsf{in }}
\newcommand{\Out}{\textsf{out }}
\newcommand{\Inout}{\textsf{inout }}

\newcommand{\encabezadoDeProc}[4]{%
	% Ponemos la palabrita problema en tt
	%  \noindent%
	{\normalfont\bfseries\ttfamily proc}%
	% Ponemos el nombre del problema
	\ %
	{\normalfont\ttfamily #2}%
	\
	% Ponemos los parametros
	(#3)%
	\ifthenelse{\equal{#4}{}}{}{%
		% Por ultimo, va el tipo del resultado
		\ : #4}
}

\newenvironment{proc}[4][res]{%
	
	% El parametro 1 (opcional) es el nombre del resultado
	% El parametro 2 es el nombre del problema
	% El parametro 3 son los parametros
	% El parametro 4 es el tipo del resultado
	% Preambulo del ambiente problema
	% Tenemos que definir los comandos requiere, asegura, modifica y aux
	\newcommand{\requiere}[2][]{%
		{\normalfont\bfseries\ttfamily requiere}%
		\ifthenelse{\equal{##1}{}}{}{\ {\normalfont\ttfamily ##1} :}\ %
		\{\ensuremath{##2}\}%
		{\normalfont\bfseries\,\par}%
	}
	\newcommand{\asegura}[2][]{%
		{\normalfont\bfseries\ttfamily asegura}%
		\ifthenelse{\equal{##1}{}}{}{\ {\normalfont\ttfamily ##1} :}\
		\{\ensuremath{##2}\}%
		{\normalfont\bfseries\,\par}%
	}
	\renewcommand{\aux}[4]{%
		{\normalfont\bfseries\ttfamily aux\ }%
		{\normalfont\ttfamily ##1}%
		\ifthenelse{\equal{##2}{}}{}{\ (##2)}\ : ##3\, = \ensuremath{##4}%
		{\normalfont\bfseries\,;\par}%
	}
	\renewcommand{\pred}[3]{%
		{\normalfont\bfseries\ttfamily pred }%
		{\normalfont\ttfamily ##1}%
		\ifthenelse{\equal{##2}{}}{}{\ (##2) }%
		\{%
		\begin{adjustwidth}{+5em}{}
			\ensuremath{##3}
		\end{adjustwidth}
		\}%
		{\normalfont\bfseries\,\par}%
	}
	
	\newcommand{\res}{#1}
	\vspace{1ex}
	\noindent
	\encabezadoDeProc{#1}{#2}{#3}{#4}
	% Abrimos la llave
	\par%
	\tocarEspacios
}
{
	% Cerramos la llave
	\vspace{1ex}
}

\newcommand{\aux}[4]{%
	{\normalfont\bfseries\ttfamily\noindent aux\ }%
	{\normalfont\ttfamily #1}%
	\ifthenelse{\equal{#2}{}}{}{\ (#2)}\ : #3\, = \ensuremath{#4}%
	{\normalfont\bfseries\,;\par}%
}

\newcommand{\pred}[3]{%
	{\normalfont\bfseries\ttfamily\noindent pred }%
	{\normalfont\ttfamily #1}%
	\ifthenelse{\equal{#2}{}}{}{\ (#2) }%
	\{%
	\begin{adjustwidth}{+2em}{}
		\ensuremath{#3}
	\end{adjustwidth}
	\}%
	{\normalfont\bfseries\,\par}%
}

% Tipos

\newcommand{\nat}{\ensuremath{\mathds{N}}}
\newcommand{\ent}{\ensuremath{\mathds{Z}}}
\newcommand{\float}{\ensuremath{\mathds{R}}}
\newcommand{\bool}{\ensuremath{\mathsf{Bool}}}
\newcommand{\cha}{\ensuremath{\mathsf{Char}}}
\newcommand{\str}{\ensuremath{\mathsf{String}}}

% Logica

\newcommand{\True}{\ensuremath{\mathrm{true}}}
\newcommand{\False}{\ensuremath{\mathrm{false}}}
\newcommand{\Then}{\ensuremath{\rightarrow}}
\newcommand{\Iff}{\ensuremath{\leftrightarrow}}
\newcommand{\implica}{\ensuremath{\longrightarrow}}
\newcommand{\IfThenElse}[3]{\ensuremath{\mathsf{if}\ #1\ \mathsf{then}\ #2\ \mathsf{else}\ #3\ \mathsf{fi}}}
\newcommand{\yLuego}{\land _L}
\newcommand{\oLuego}{\lor _L}
\newcommand{\implicaLuego}{\implica _L}

\newcommand{\cuantificador}[5]{%
	\ensuremath{(#2 #3: #4)\ (%
		\ifthenelse{\equal{#1}{unalinea}}{
			#5
		}{
			$ % exiting math mode
			\begin{adjustwidth}{+2em}{}
				$#5$%
			\end{adjustwidth}%
			$ % entering math mode
		}
		)}
}

\newcommand{\existe}[4][]{%
	\cuantificador{#1}{\exists}{#2}{#3}{#4}
}
\newcommand{\paraTodo}[4][]{%
	\cuantificador{#1}{\forall}{#2}{#3}{#4}
}

%listas

\newcommand{\TLista}[1]{\ensuremath{seq \langle #1\rangle}}
\newcommand{\lvacia}{\ensuremath{[\ ]}}
\newcommand{\lv}{\ensuremath{[\ ]}}
\newcommand{\longitud}[1]{\ensuremath{|#1|}}
\newcommand{\cons}[1]{\ensuremath{\mathsf{addFirst}}(#1)}
\newcommand{\indice}[1]{\ensuremath{\mathsf{indice}}(#1)}
\newcommand{\conc}[1]{\ensuremath{\mathsf{concat}}(#1)}
\newcommand{\cab}[1]{\ensuremath{\mathsf{head}}(#1)}
\newcommand{\cola}[1]{\ensuremath{\mathsf{tail}}(#1)}
\newcommand{\sub}[1]{\ensuremath{\mathsf{subseq}}(#1)}
\newcommand{\en}[1]{\ensuremath{\mathsf{en}}(#1)}
\newcommand{\cuenta}[2]{\mathsf{cuenta}\ensuremath{(#1, #2)}}
\newcommand{\suma}[1]{\mathsf{suma}(#1)}
\newcommand{\twodots}{\ensuremath{\mathrm{..}}}
\newcommand{\masmas}{\ensuremath{++}}
\newcommand{\matriz}[1]{\TLista{\TLista{#1}}}
\newcommand{\seqchar}{\TLista{\cha}}

\renewcommand{\lstlistingname}{Código}
\lstset{% general command to set parameter(s)
	language=Java,
	morekeywords={endif, endwhile, skip},
	basewidth={0.47em,0.40em},
	columns=fixed, fontadjust, resetmargins, xrightmargin=5pt, xleftmargin=15pt,
	flexiblecolumns=false, tabsize=4, breaklines, breakatwhitespace=false, extendedchars=true,
	numbers=left, numberstyle=\tiny, stepnumber=1, numbersep=9pt,
	frame=l, framesep=3pt,
	captionpos=b,
}

\newcommand{\notimplies}{\;\not\!\!\!\implies}
\title{Álgebra I}
\author{Tomás Agustín Hernández}
\date{}

\begin{document}
\maketitle

\begin{figure}[b]
    \centering
    \begin{tikzpicture}[remember picture,overlay]
        \node[anchor=south east, inner sep=0pt, xshift=-1cm, yshift=2cm] at (current page.south east) {
            \begin{minipage}[b]{0.5\textwidth}
                \includegraphics[width=\linewidth]{logo_uba.jpg}
                \label{fig:bottom}
            \end{minipage}
        };
    \end{tikzpicture}
\end{figure}

\newpage
\section*{Conjuntos}
Los conjuntos almacenan elementos, no se consideran repetidos y responde a la pregunta de $"$¿está el elemento?$"$, esto último quiere decir que no tenemos forma de tomar un elemento sino predicar acerca de si está o no.
\subsection*{Pertenecencia a un Conjunto}
Si consideramos cualquier elemento $x$, decimos que está en un conjunto A si \textbf{x pertenece a A}. \\
La pertenencia de un elemento a un conjunto la denotamos como: $x \in A$ \\
\textbf{Importante}: La relación está dada por $Elemento \ en \ Conjunto$ \\
Véase \hyperref[subsec:pertenecencia_conjuntos]{\underline{ánexo}} para ejemplos más didácticos.
\subsection*{Inclusión a un Conjunto}
Sean A y D conjuntos cualesquiera.
Decimos que D es un subconjunto de A sí y solo sí todos los elementos de D están en A. \\
La inclusión en un conjunto la denotamos como $D \subseteq A$ \\
Es posible leer el símbolo $ \subseteq $ de tres maneras: 
\begin{itemize}
    \item "D es un subconjunto de A"
    \item "D está incluido en A"
    \item "D está contenido en A"
\end{itemize}
Los subconjuntos posibles no salen más que haciendo combinaciones con sus elementos, es decir, agruparlos de diferentes formas. \\
Véase \hyperref[subsec:inclusion_conjuntos]{\underline{ánexo}} para ejemplos más didácticos.
\subsection*{Cardinal de un Conjunto}
Sea A un conjunto, el cardinal de un conjunto indica la cantidad de elementos en el conjunto. \\
Se denota como: $\#A$
\subsection*{Cantidad de Subconjuntos posibles dado un Conjunto}
Sea un conjunto A, la cantidad de subconjuntos D para el conjunto A es: $2^{\#A}$
\subsection*{Elemento Vacío}
Se representa con el símbolo de $\emptyset$. El elemento vacío está \textbf{incluido} en todos los conjuntos. \\
\textbf{Importante}: El elemento vacío NO pertenece a todos los conjuntos sino que está incluido en todos.
\subsection*{Cuantificadores}
Nos permiten predicar acerca de los elementos de un conjunto dado. 
\begin{itemize}
    \item $\forall$ x: Para todo x.
    \begin{itemize}
        \item Para que sea verdadero todos deben cumplir la condición dada.
        \item Es falso si existe un caso en que no se cumple.
    \end{itemize}
    \item $\exists$ x: Existe un x
    \begin{itemize}
        \item Para que sea verdadero alcanza con encontrar un caso verdadero.
        \item Es falso si no hay ningun caso que cumpla la condición
    \end{itemize}
\end{itemize}
\textbf{Importante}: El símbolo de $:$ o $\symbol{92}$ significa "tal que" \\
Véase \hyperref[subsec:cuantificadores]{\underline{ánexo}} para ejemplos más didácticos.
\section*{Operaciones entre Conjuntos}
Sean A y B conjuntos cualesquiera. La cantidad de filas que tendrá una tabla de verdad es: \textbf{$2^{cantVariables}$} \\
\textbf{Importante}: Las operaciones entre conjuntos que vamos a ver están relacionadas con la lógica proposicional.
\newpage
\subsection*{Unión $(A \cup B)$}
Es exactamente igual como en la lógica proposicional. La unión es un $o$ lógico. En el conjunto resultante quedan los elementos de A y B. \\

\begin{table}[h!]
    \centering
    \begin{tabular}{|c | c | c|}
    \hline
    \textbf{A} & \textbf{B} & \textbf{$A \cup B$} \\[0.1cm]
    \hline
    V & V & V \\
    V & F & V \\
    F & V & V \\
    F & F & F \\
    \hline
    \end{tabular}
    \caption{Unión de conjuntos}
\end{table} 
Cada fila se puede generalizar para un x cualquiera en las operaciones lógicas. \\
Ej.: Si $x \in A \land x \in B$ entonces $ x \in A \cup B$ esto claramente nos dice que estamos en el caso de la fila 1. \\
Ej.: Si $x \notin A \land x \in B$ entonces $ x \in A \cup B$ esto claramente nos dice que estamos en el caso de la fila 3.
\subsection*{Intersección $(A \cap B)$}
Es exactamente igual como en la lógica proposicional. La intersección es un $"y"$ lógico. En el conjunto resultante quedan los elementos que están tanto en A y en B.
\begin{table}[h!]
    \centering
    \begin{tabular}{|c | c | c|}
    \hline
    \textbf{A} & \textbf{B} & \textbf{$A \cap B$} \\[0.1cm]
    \hline
    V & V & V \\
    V & F & F \\
    F & V & F \\
    F & F & F \\
    \hline
    \end{tabular}
    \caption{Intersección de conjuntos}
\end{table} \\
Cada fila se puede generalizar para un x cualquiera en las operacines lógicas. \\
Ej.: Si $x \in A \land x \in B$ entonces $ x \in A \cap B$ esto claramente nos dice que estamos en el caso de la fila 1. \\
Ej.: Si $x \notin A \land x \in B$ entonces $ x \notin A \cap B$ esto claramente nos dice que estamos en el caso de la fila 3.
\subsection*{Complemento $(A \cap B)$}
En la lógica proposicional, el complemento es la negación. Lo que está en un conjunto universal V pero no en el conjunto.
\begin{table}[h!]
    \centering
    \begin{tabular}{|c | c|}
    \hline
    \textbf{A} & \textbf{$\neg A$} \\[0.1cm]
    \hline
    V & F \\
    V & F  \\
    F & V  \\
    F & V \\
    \hline
    \end{tabular}
    \caption{Complemento en Conjuntos}
\end{table} \\
Cada fila se puede generalizar para un x cualquiera en las operaciones lógicas. \\
Ej.: Si $x \in A$ entonces termina siendo $ x \notin A$ esto claramente nos dice que estamos en el caso de la fila 1. \\

Sea $A=\{1, 2\}, B=\{3, 4, 5\}, C=\{8, 9\}, V=\{A, B, C\} \implies A^{c} = \{3, 4, 5, 8, 9\}$ \\
\textbf{Importante}: Nótese que siempre se hace el complemento en base a los elementos que hay en el universo y se excluyen algunos. En este caso, del universo V nos quedamos con los que NO están en A.
\newpage
\subsection*{Diferencia $(A-B )$}
Esta operación es conocida también de la siguiente manera $A\symbol{92}B$.
Es una equivalencia de $A \cap B^{c}$. Representa lo que está en A pero no en B. Si se lo quisiera representar en la tabla de verdad, debe representar la equivalencia.
\begin{table}[h!]
    \centering
    \begin{tabular}{|c | c | c | c|}
    \hline
    \textbf{A} & \textbf{B} & \textbf{$B^{c}$} & \textbf{$A \cap B^{c}$} \\[0.1cm]
    \hline
    V & V & F & F\\
    V & F & V & V\\
    F & V & F & F\\
    F & F & V & F\\
    \hline
    \end{tabular}
    \caption{Diferencia de conjuntos}
\end{table} \\
\subsection*{Diferencia Simétrica $(A \triangle B)$}
Equivalente al XOR($\veebar$) u $o$ excluyente en la lógica proposicional. \\
Es una equivalencia de $(A - B) \cup (B - A)$ y $(A\cup B) - (A \cap B)$. Representa lo que está en A o en B pero no en ambos. 
\begin{table}[h!]
    \centering
    \begin{tabular}{|c | c | c | c | c|}
    \hline
    \textbf{A} & \textbf{B} & \textbf{$A \veebar B$} & \textcolor{blue}{$(A - B) \cup (B - A)$} & \textcolor{blue}{$(A\cup B) - (A \cap B)$} \\[0.1cm]
    \hline
    V & V & F & F & F \\
    V & F & V & V & V \\
    F & V & V & V & V \\
    F & F & F & V & V \\
    \hline
    \end{tabular}
    \caption{Diferencia Simétrica en conjuntos}
\end{table} \\
\textbf{Nota}: Las columnas en azul son equivalencias a la operación $\veebar$ y son útiles a la hora de demostrar. \\
Ej.: Si $x \in A \land x \in B$ entonces $ x \veebar B = F$ esto claramente nos dice que estamos en el caso de la fila 1.
Ej.: Si $x \in A \land x \notin B$ entonces $ x \veebar B = V$ esto claramente nos dice que estamos en el caso de la fila 2. \\
\subsection*{Inclusión $(A \subseteq B)$}
Representa el $\implies$ de la lógica proposicional. Recordemos que la inclusión es verdadera si todos los elementos de A están en B siendo A y B conjuntos cualesquiera. \\
Es lo que vamos a utilizar para demostrar, y es importante que se lo entienda bien. 
\begin{table}[h!]
    \centering
    \begin{tabular}{|c | c | c|}
    \hline
    \textbf{A} & \textbf{B} & \textbf{$A \implies B$} \\[0.1cm]
    \hline
    V & V & V \\
    V & F & F \\
    F & V & V \\
    F & F & V \\
    \hline
    \end{tabular}
    \caption{Inclusión de conjuntos}
\end{table} 
\textbf{Tips}: 
\begin{itemize}
    \item El único caso que nos importa es que si el antecedente es verdadero, hay que ver que el consecuente NO sea falso. En las demostraciones asumimos que vale el antecedente y tenemos que ver si hace verdadero al consecuente.
    \item Si no se cumple el antecedente, el consecuente es siempre verdadero.
\end{itemize}
Cada fila se puede generalizar para un x cualquiera en las operacines lógicas. \\
Ej.: Sea $A = \{1, 2, 3\} \ B = \{10, 40\} \ x = 100$ ¿Se cumple que $ x \in A \implies x \in B$? ¿100 está en A? No, y al ser una implicación si el antecedente no se cumple, queda toda la proposición verdadera. Luego, sí, se cumple que $ x \in A \implies x \in B$. Esto claramente nos dice que estamos en el caso de la fila 3. \\ 
Ej.: Sea $A = \{1, 2, 3\} \ B = \{10, 40\} \ x = 3$ ¿Se cumple que $ x \in A \implies x \in B$? ¿3 está en A? Sí. Entonces esto hace al antecedente verdadero ¿me basta para decir que la proposición es verdadera? No. Primero debo ver qué pasa con el consecuente. ¿Es cierto que 3 está en B? No. Entonces como el antecedente es verdadero y el consecuente es falso, la proposición es falsa. Luego, no, no se cumple que $ x \in A \implies x \in B$. Esto claramente nos dice que estamos en el caso de la fila 2.
\section*{Anexo}
\subsection*{Pertenencia en Conjuntos}
\label{subsec:pertenecencia_conjuntos}
Sea A el conjunto: $\{1, 2, \{C, B\}, F, \{10, 15\}\}$
\begin{itemize}
    \item $ 1 \in A, 2 \in A, F \in A $
    \item $ C \notin A, B \notin A $
    \item $  \{C, B\}, \{10, 15\} \in A $
\end{itemize}
¿Por qué $C \notin A$? Pues C no es un elemento de A.\\ Notar que C es parte del elemento $\{C, B\}$ en A, pero C no es un elemento independiente.
\subsection*{Inclusión en Conjuntos}
\label{subsec:inclusion_conjuntos}
\textbf{Ex. 1}: Sea $A = \{1, 2, 3\} \ y \ D = \{1, 3\}$. ¿Es D un subconjunto de A? \\
Sí, lo es pues $1 \in A$ y $ 3 \in A$ \\
\textbf{Ex. 2}: Sea $A = \{1, \{1, 4\}, 3, 10\}$
\begin{itemize}
    \item $ \{1, 4\} \nsubseteq A $ pues no existen 1 y 4 como elementos en A
    \item $ \{1, 4\} \in A $ pues $\{1, 4\} es un elemento de A$
    \item $\{1, 3\} \subseteq A$ pues $ 1 \in A, 3 \in A$, lo mismo sucede con $ \{1, 10\} \ o \ \{3, 10\}$
\end{itemize}
\subsection*{Cuantificadores}
\label{subsec:cuantificadores}
\textbf{Ex. 1}: $A = \{2, 4, 6, 8\}$ \\
Algunos ejemplos utilizando cuantificadores 
\begin{itemize}
    \item $\forall x \in A \ \symbol{92} \ x \% 2 = 0$ (Todos pares en A)
    \item $\neg \ \exists x \in A \ \symbol{92} \ x \% 2 \neq 0$ (No existe ningún impar en A) 
    \item $ \exists x \in A \ \symbol{92} x = 4$ (Existe un elemento en A que es exactamente 4)
\end{itemize}

\end{document} 

